%% =========================================================
%% Majorization and Birkhoff's theorem
%% Doubly-stochastic half.
%% =========================================================

\section{Majorization}

Majorization is a preorder on real vectors that captures when one
probability distribution is more ``mixed'' than another.
Its convex-geometric foundation is Birkhoff's theorem.

\begin{theorem}[Birkhoff]
    \label{thm:birkhoff_doubly_stochastic}
    \lean{TNLean.wolf_birkhoff_doublyStochastic_convexHull}
    \lean{TNLean.wolf_birkhoff_extremePoints_doublyStochastic}
    \leanok
    The set of doubly stochastic matrices in $M_d(\R)$ is the convex
    hull of the $d\times d$ permutation matrices, and its extreme
    points are exactly the permutation matrices.
    See \cite[Chapter~8, lines~263--266]{Wolf2012Quantum}.
\end{theorem}
% The doubly-substochastic half of Birkhoff's theorem (extreme
% points of the substochastic matrices) is not formalized; #3959.

\begin{proof}\leanok
    Every doubly stochastic matrix is a convex combination of
    permutation matrices, and every convex combination of permutation
    matrices is doubly stochastic (the first identity).  Since the
    permutation matrices themselves are extreme points of the doubly
    stochastic matrices and no other extreme points exist, the second
    identity follows.
\end{proof}
\endinput
